% C.3bd. Final corrected proof; full root and real_weil audits passed.
\paragraph{C.3bd. Top Chow groups of isolated surface and threefold
germs over \(\overline{\mathbb F}_p\) (Proved).}
Let \(k=\overline{\mathbb F}_p\). Let \(R\) be the strict
henselization at a closed \(k\)-point of a normal local ring
essentially of finite type over \(k\), and suppose that
\[
                  2\leq N=\dim R\leq3,\qquad
                  X=\operatorname{Spec}R\setminus\{\mathfrak m_R\}
                  \quad\hbox{is regular}.
\]
Then
\[
                         \mathrm{CH}^{N-1}(X)_{\mathbb Q}=0,
                 \qquad B^{N-1}(X)_{\mathbb Q}=0.
\]
In dimension two the regular-puncture hypothesis follows from
normality. We retain it explicitly in dimension three. The assertion
holds in every characteristic \(p>0\), including \(2\) and \(3\).
It is a low-dimensional consequence of C.3bc, not an assertion of
resolution in higher dimension.

\begin{proof}
\emph{A projective modification preserving the puncture.}
Put \(S=\operatorname{Spec}R\). If \(R\) is regular, its blow-up at
the closed point has exceptional divisor \(\mathbb P^{N-1}_k\)
and relatively ample invertible sheaf \(\mathcal O(-E)\).
We can therefore apply C.3bc directly.

Suppose now that the closed point is singular.
The resolution theorem of
\citet[Theorem 1.1, current author version pp.~3--4]
{frC3bdCossartPiltant2019}, applied to the affine scheme \(S\),
gives a projective resolution \(\pi_1:Z_1\to S\) that is an
isomorphism over \(X\). We make the ample-exceptional
construction explicit, so projectivity and preservation of
\(X\) are not inferred from separate existence assertions.
Choose a \(\pi_1\)-ample invertible sheaf \(H_1\) and put
\(\mathcal L_1=(\pi_1)_*H_1\), a coherent generically rank-one
sheaf, invertible on \(X\). Set
\(\mathcal L_2=\mathcal Hom(\mathcal L_1,\mathcal O_S)\), and
let \(\pi_2:Z_2\to S\) be the Proj construction in
\citet[Lemma 8, author pp.~2--3]{frC3bdKollarWitaszek2026}.
It is the identity on \(X\).

Let \(M\) be the torsion-free quotient of
\(\pi_1^*\mathcal L_2\). Since \(Z_1\) is regular, its
rank-one reflexive hull \(M^{**}\) is invertible. Thus
\[
                    J=M\otimes(M^{**})^{-1}
                              \subset\mathcal O_{Z_1}
\]
is a coherent ideal, equal to the unit ideal on \(X\) and
at every codimension-one point. If \(N=3\), apply
\citet[Proposition 4.4, author p.~141]
{frC3bdCossartPiltant2019}: a finite sequence of blow-ups
in regular centres, supported in the nonprincipal locus of
the successive ideal transforms, makes \(J\) invertible.
If \(N=2\), the support of \(\mathcal O_{Z_1}/J\) is a finite
set of closed points over \(\mathfrak m_R\), and the same
conclusion follows from
\citet[Lemma 54.4.1, Tag 0AHH]{frC3bdStacksPrincipalization}.
In either case we obtain a regular scheme \(Z_0\), projective
over \(Z_1\), unchanged over \(X\).
The pulled-back sheaf \(\mathcal L_2\) has an invertible
quotient there: it is \(J\mathcal O_{Z_0}\) tensored with the
pullback of \(M^{**}\). The universal property of the Proj
construction therefore gives \(Z_0\to Z_2\).
We can now apply Koll\'ar--Witaszek's Lemma~8 directly to
the two maps \(Z_0\to Z_1,Z_2\). It gives a regular
projective modification
\[
                              g_0:Z_0\longrightarrow S
\]
which is an isomorphism over \(X\), together with a
\(g_0\)-ample Cartier divisor \(F_0\) supported over
\(\mathfrak m_R\).
This construction does not require a further projective
resolution of a nonaffine fibre-product component. Every
additional blow-up has its centre outside \(X\).

\emph{Making the exceptional fibre SNC without losing ampleness.}
The entire reduced closed fibre
\((g_0^{-1}(\mathfrak m_R))_{\mathrm{red}}\) has dimension at most
\(N-1\leq2\). Apply the embedded-resolution result
\citet[Corollary 0.4 and its proof, author p.~5]
{frC3bdCossartJannsenSaito2013} to this reduced closed subset
inside the regular excellent scheme \(Z_0\).
It supplies a projective morphism
\[
                              h:Z\longrightarrow Z_0
\]
with regular \(Z\), an isomorphism away from that subset, whose
reduced inverse image is exactly the full reduced closed fibre,
an SNC divisor \(E\).
Consequently \(f=g_0h\) is still an isomorphism over \(X\).
The scheme \(Z_0\) is regular, hence locally factorial.
Koll\'ar--Witaszek's Lemma~6, author p.~2, gives an
\(h\)-ample \(h\)-exceptional Cartier divisor \(G\).
After taking a positive multiple of \(G\), the usual relative
ampleness composition gives, for \(m\) sufficiently large,
\[
                          F=h^*(mF_0)+G
\]
an \(f\)-ample Cartier divisor. This is also the ampleness
calculation used in the proof of their Lemma~8, author pp.~2--3.
Every component of \(F\) is supported over the closed point.

\emph{The sign and full support are forced by normality.}
We verify the positive-coefficient hypothesis of C.3bc rather
than imposing a sign convention on the phrase
\emph{ample exceptional divisor}.
Since \(S\) is affine, \(f\)-ampleness and
\citet[Proposition 28.27.13(7), Tag 01Q3]
{frC3bdStacksGeneration} show that
\(\mathcal O_Z(nF)\) is generated by its global sections for
all sufficiently large \(n\).
Using the canonical rational-function trivialization of this
divisor sheaf on \(X\), set
\[
                         I=f_*\mathcal O_Z(nF).
\]
Normal Hartogs gives
\(\Gamma(X,\mathcal O_X)=R\)
\citep[Remark 31.17.2, Tag 0BCS]{frC3bdStacksHartogs}.
It therefore identifies \(I\) with a coherent ideal of \(R\),
equal to the unit ideal away from \(\mathfrak m_R\).
The evaluation surjection is the natural inclusion of fractional
ideal sheaves in the common function field, so
\[
                      I\mathcal O_Z=\mathcal O_Z(nF)
                                     \subset\mathcal O_Z.
\]
Thus \(-nF\) is effective. If \(I=R\), this equality makes
\(\mathcal O_Z(nF)\) trivial; it cannot be relatively ample on
a positive-dimensional projective fibre. The resulting
quasi-finite projective birational morphism to normal \(S\)
would be an isomorphism, contrary to the singularity of its
closed point. Hence \(I\) is a proper
\(\mathfrak m_R\)-primary ideal. Its pullback vanishes on
exactly the entire closed fibre. Consequently
\[
                    -nF=\sum_i a_iE_i,\qquad a_i>0,
\]
and \(nF|_E\) is ample. These are exactly the exceptional
divisor hypotheses of C.3bc. That theorem now gives
\(\mathrm{CH}^{N-1}(X)_{\mathbb Q}=0\).
Its ordinary quotient is zero as well.
\end{proof}
